% Supplementary Note 5: Study 5, the label's effect at each stated editing time. Numbers through macros only.
\suppnote{Study 5: the label's effect at each stated editing time}

This note gives, cell by cell, the exploratory breakdown behind Study~5's one significant registered confirmatory test: the AI label's creativity penalty differed across stated editing times. The main text reports that test in its Results section on stated editing hours.

\paragraph*{Design.}
In Study~5, \SFiveN{} participants each saw six original images, all generated with an AI image generator. With each original came three second works that Alex was said to have made from it with his own design and sketching tools: None, the original shown again unchanged; Slight, a light edit; and Dramatic, a heavy one. Every participant saw all three levels. Two factors varied between participants, so each person saw one level of each for the whole session. The label said that Alex had made the originals either with design and sketching tools (the Human label) or by prompting an AI system (the AI label); because every original was AI-generated, the Human label was a cover story. The stated editing time said that Alex spent under 1 hour, 10 hours or 100 hours making each second work. The six groups held \SFiveNHumanUnderOneHour, \SFiveNHumanTenHours{} and \SFiveNHumanHundredHours{} participants under the Human label and \SFiveNAIUnderOneHour, \SFiveNAITenHours{} and \SFiveNAIHundredHours{} under the AI label, at under 1, 10 and 100 hours (descriptive). The sentence that stated the editing time was the same at every level, so at None the description participants read said that Alex had spent the stated hours on an unchanged copy. A time for the original was also shown on every trial but was not recorded, so it enters no analysis. Participants rated each second work's transformativeness, essence change and creativity on 1-to-5 scales. As in Studies 2 and 4, the model's predicted means rose with the level of modification; on creativity under the Human label they were \SFiveCreatMeanHumanNone, \SFiveCreatMeanHumanSlight{} and \SFiveCreatMeanHumanDramatic{} at None, Slight and Dramatic (descriptive). Study~5's registration was filed on 8 October 2026, after data collection but before any of its ratings had been averaged or compared by label.

\paragraph*{The penalty.}
Each estimate below is the label's penalty: the predicted mean rating under the Human label minus the predicted mean under the AI label, in points on the 1-to-5 scale, from Study~5's registered linear mixed model, a regression that gives each participant and each image its own baseline. A positive penalty means that the AI label lowered the rating, and a negative one that it raised it. A $p$-value is the chance of a difference at least this large if there were no effect, and a result is called significant when $p$ is below .05.

\paragraph*{The registered test.}
Study~5's registration named \SFiveNConfirmatory{} confirmatory tests, fixed with their decision rule before the analysis: the effect of editing time and the label-by-editing-time interaction, which asks whether the label's penalty differs across stated editing times, on each of the three ratings. They were judged together after Holm's correction, which raises each $p$-value in a set of tests so that the chance of calling any result in the set significant by mistake stays at .05. Only \SFiveNConfirmatoryHolmSignificant{} of them passed: the label's creativity penalty differed across stated editing times ($F$(\SFiveCreatSourceByEditTimeDfNum,~\SFiveCreatSourceByEditTimeDfDen)~=~\SFiveCreatSourceByEditTimeF, uncorrected $p$~\SFiveCreatSourceByEditTimeP, Holm $p$~\SFiveCreatSourceByEditTimePHolm). Two checks that the registration named as secondary agreed. An ordinal model, which treats the five answers as ordered categories rather than evenly spaced numbers, found the same interaction ($\chi^2$~=~\SFiveCreatSourceByEditTimeOrdChisq, Holm $p$~\SFiveCreatSourceByEditTimeOrdPHolm) and nothing else among the six tests. The registered second analysis, with every item rescaled from 0 to 1, gave the same verdicts ($F$~=~\SFiveCreatSourceByEditTimeRescaledF, Holm $p$~\SFiveCreatSourceByEditTimeRescaledPHolm). The interaction was not significant on the other two measures (Holm $p$~\SFiveTransSourceByEditTimePHolm{} on transformativeness and $p$~\SFiveEssSourceByEditTimePHolm{} on essence change). Averaged over the two labels, stated editing hours did not detectably change ratings on any measure (registered confirmatory; reported in the main text). The interaction concerns the penalty, not the mean rating, so the two results do not conflict. That null cannot be set against Study~3's effect of stated time on originals as a finding, because in Study~3 each participant saw every stated time and in Study~5 only one.

\begin{figure}[tbp]
\centering
\makebox[\textwidth][c]{\includegraphics[width=183mm]{figS_study5_editing_time}}
\caption{\textbf{Exploratory: Study~5's creativity gap at each stated editing time.} Each square is the gap in creativity ratings of the second work, the Human-label mean minus the AI-label mean, from Study~5's registered statistical model fitted again to give a gap at each editing time and level, with its 95\% confidence interval, not corrected for the nine comparisons. Right of zero, the AI label lowered the rating. Editing time varied between participants, so each panel compares two separate groups of \SFiveNCellMin{} to \SFiveNCellMax{} people, one per label. None is the original shown again unchanged. The registered test behind this breakdown, whether the label's creativity gap differed across editing times, was significant after Holm's correction (Holm $p$~\SFiveCreatSourceByEditTimePHolm). The breakdown was specified on 9 October 2026, after the registered results were known, and it changes no registered verdict.}
\label{figS:editingtime}
\end{figure}

\paragraph*{The shape of the interaction.}
The registration sketched one shape: if editing counts as Alex's own contribution, the penalty should shrink as the hours rise. The data took another. An exploratory breakdown, specified on 9 October 2026 after the registered results were known, estimated the creativity penalty at each editing time and level from Study~5's registered model, fitted again for this purpose (Supplementary Fig.~\ref{figS:editingtime}). It carries no verdict, and its $p$-values are not corrected for the nine comparisons. At under 1 hour, the penalty was small at every level: \XSTimeCheckSFiveCreatPenaltyUnderOneHourNoneEst, \XSTimeCheckSFiveCreatPenaltyUnderOneHourSlightEst{} and \XSTimeCheckSFiveCreatPenaltyUnderOneHourDramaticEst{} points at None, Slight and Dramatic. At 10 hours it was largest at every level: \XSTimeCheckSFiveCreatPenaltyTenHoursNoneEst, \XSTimeCheckSFiveCreatPenaltyTenHoursSlightEst{} and \XSTimeCheckSFiveCreatPenaltyTenHoursDramaticEst{} (uncorrected $p$~\XSTimeCheckSFiveCreatPenaltyTenHoursSlightP{} at Slight and $p$~\XSTimeCheckSFiveCreatPenaltyTenHoursDramaticP{} at Dramatic). At Slight and Dramatic these are about as large as Study~1's creativity penalty on originals, \SOneHOneCreatDiff{} points, a comparison across studies rather than a test. At 100 hours the penalty was small at None and Slight (\XSTimeCheckSFiveCreatPenaltyHundredHoursNoneEst{} and \XSTimeCheckSFiveCreatPenaltyHundredHoursSlightEst) but reversed at Dramatic (\XSTimeCheckSFiveCreatPenaltyHundredHoursDramaticEst, uncorrected $p$~\XSTimeCheckSFiveCreatPenaltyHundredHoursDramaticP). The model term that lets this editing-time pattern differ across levels was also significant on creativity ($F$(\SFiveCreatSourceByEditTimeByModDfNum,~\SFiveCreatSourceByEditTimeByModDfDen)~=~\SFiveCreatSourceByEditTimeByModF, uncorrected $p$~\SFiveCreatSourceByEditTimeByModP); it is exploratory, because the registration names no test of it. The same breakdown on the other two measures, whose registered interactions were not significant, also put the largest penalty in the 10-hour groups at every level; at Slight and Dramatic it was \XSTimeCheckSFiveTransPenaltyTenHoursSlightEst{} and \XSTimeCheckSFiveTransPenaltyTenHoursDramaticEst{} on transformativeness and \XSTimeCheckSFiveEssPenaltyTenHoursSlightEst{} and \XSTimeCheckSFiveEssPenaltyTenHoursDramaticEst{} on essence change.

\paragraph*{What the breakdown means for the Slight results.}
Study~5 registered the penalty at each level as secondary tests, judged at .05 without correction. Each was read with an equivalence test, which asks whether the penalty's 90\% confidence interval, the range of values the data do not rule out at the 90\% level, lies wholly inside $\pm$\SourceBound{} points, half of Study~1's creativity penalty on originals. At Slight the penalty was significant and not shown inside that bound on all three measures; on creativity it was \SFiveCreatPenaltySlightEst{} points ($p$~\SFiveCreatPenaltySlightP; Supplementary Table~\ref{tab:supp-labeleffects}). That estimate averages over the three editing times; taken one time at a time, the Slight creativity penalty was \XSTimeCheckSFiveCreatPenaltyUnderOneHourSlightEst, \XSTimeCheckSFiveCreatPenaltyTenHoursSlightEst{} and \XSTimeCheckSFiveCreatPenaltyHundredHoursSlightEst{} at under 1, 10 and 100 hours (exploratory), so the residue at Slight, what the label carried into ratings of slight edits, rests largely on the 10-hour groups.

\paragraph*{Readings the data cannot separate.}
Each combination of label and editing time was a separate group of \SFiveNCellMin{} to \SFiveNCellMax{} people. One unusual 10-hour group, AI-label raters rating low or Human-label raters rating high, would produce this pattern on all three measures, and these data cannot rule that out; if so, it also inflates Study~5's Slight estimates in Supplementary Table~\ref{tab:supp-labeleffects}. The corrected $p$ sits just under .05, and none of the accounts considered in the paper predicts a peak at 10 hours, so the interaction may be a chance result that survived the correction.
